% =====================================================================
%  BÁO CÁO PROJECT II – ĐỀ TÀI 1A
%  Biên dịch: latexmk -xelatex main.tex   (hoặc recipe "latexmk (XeLaTeX)" trong LaTeX Workshop)
% =====================================================================
\documentclass[13pt,a4paper,oneside]{extreport}

% ---------- Trang & phông chữ (quy cách đồ án: Times New Roman 13, lề trái 3 cm)
\usepackage[left=3cm,right=2cm,top=2.5cm,bottom=2.5cm,headheight=17pt]{geometry}
\usepackage{fontspec}
\usepackage{polyglossia}
\setdefaultlanguage{vietnamese}
\setotherlanguage{english}
\setmainfont{Times New Roman}
\setsansfont{Arial}
\setmonofont[Scale=0.84]{Consolas}
\usepackage{setspace}
\setstretch{1.3}
\setlength{\parskip}{4pt}
\setlength{\parindent}{0.8cm}

% ---------- Gói hỗ trợ
\usepackage{graphicx}
\usepackage{float}
\usepackage[labelsep=period,font=small,labelfont=bf,skip=6pt]{caption}
\usepackage[table]{xcolor}
\usepackage{booktabs,tabularx,longtable,array,multirow,makecell}
\usepackage{xltabular}
\usepackage{enumitem}
\usepackage{fancyhdr}
\usepackage{titlesec}
\usepackage{tikz}
\usepackage[most]{tcolorbox}
\usepackage{listings}
\usepackage{fancyvrb}
% Khối văn bản thuần (cây thư mục, lệnh, kết quả chạy) – fancyvrb giữ nguyên ký tự Unicode
\DefineVerbatimEnvironment{PlainCode}{Verbatim}{fontsize=\footnotesize,frame=single,framesep=6pt,rulecolor=\color{black!15},baselinestretch=1.05}
\usepackage{amsmath,amssymb}
\usepackage{tocloft}
\usepackage[hidelinks,unicode]{hyperref}

\graphicspath{{images/}{images/screens/}{../docs/diagrams/png/}}

% ---------- Màu nhận diện
\definecolor{hustred}{HTML}{B3261E}
\definecolor{navy}{HTML}{14213D}
\definecolor{blueA}{HTML}{1F4E8C}
\definecolor{lightgray}{HTML}{F4F6FA}
\definecolor{codebg}{HTML}{F7F8FB}
\definecolor{codecomment}{HTML}{6A7A90}
\definecolor{codekw}{HTML}{B3261E}
\definecolor{codestr}{HTML}{1E7B45}

% ---------- Tiêu đề chương / mục
\titleformat{\chapter}[display]{\centering\bfseries\color{navy}}{\large CHƯƠNG \thechapter}{6pt}{\LARGE\MakeUppercase}
\titlespacing*{\chapter}{0pt}{-10pt}{22pt}
\titleformat{\section}{\large\bfseries\color{navy}}{\thesection}{0.6em}{}
\titleformat{\subsection}{\normalsize\bfseries\color{navy}}{\thesubsection}{0.6em}{}
\titleformat{\subsubsection}{\normalsize\bfseries\itshape}{\thesubsubsection}{0.6em}{}
\setcounter{secnumdepth}{3}
\setcounter{tocdepth}{2}
\renewcommand{\cftchapfont}{\bfseries}
\renewcommand{\cftchappagefont}{\bfseries}
\renewcommand{\cftchappresnum}{Chương }
\setlength{\cftchapnumwidth}{5.5em}
\renewcommand{\cftchapleader}{\cftdotfill{\cftdotsep}}
\renewcommand{\cftfigpresnum}{Hình }
\renewcommand{\cfttabpresnum}{Bảng }
\setlength{\cftfignumwidth}{5em}
\setlength{\cfttabnumwidth}{5em}

% ---------- Đầu trang / chân trang
\pagestyle{fancy}
\fancyhf{}
\fancyhead[L]{\small\color{navy}Project II – Đề tài 1A}
\fancyhead[R]{\small\color{navy}\nouppercase{\leftmark}}
\fancyfoot[C]{\small\thepage}
\renewcommand{\headrulewidth}{0.4pt}
\renewcommand{\chaptermark}[1]{\markboth{Chương \thechapter. #1}{}}
\fancypagestyle{plain}{\fancyhf{}\fancyfoot[C]{\small\thepage}\renewcommand{\headrulewidth}{0pt}}

% ---------- Danh sách, bảng
\setlist{itemsep=2pt,topsep=4pt,leftmargin=1.4em}
\renewcommand{\arraystretch}{1.25}
\newcolumntype{L}[1]{>{\raggedright\arraybackslash}p{#1}}
\newcolumntype{Y}{>{\raggedright\arraybackslash}X}
\newcommand{\hd}[1]{\textbf{\color{white}#1}}
\newcommand{\headrow}{\rowcolor{navy}}
% Khung lưới cho bảng: đường kẻ màu xám đậm, đủ rõ khi in
\setlength{\arrayrulewidth}{0.6pt}
\arrayrulecolor{black!55}

% ---------- Hình: \fig[chiều rộng]{file}{chú thích}{nhãn}
\newcommand{\fig}[4][0.95\linewidth]{%
  \begin{figure}[H]\centering
    \includegraphics[width=#1,height=0.78\textheight,keepaspectratio]{#2}
    \caption{#3}\label{#4}
  \end{figure}}

% ---------- Hộp ghi chú
\newtcolorbox{notebox}[1][]{colback=lightgray,colframe=blueA,boxrule=0.6pt,arc=2pt,left=8pt,right=8pt,top=6pt,bottom=6pt,fonttitle=\bfseries,title={#1}}
\newtcolorbox{rulebox}[1][]{colback=hustred!5,colframe=hustred,boxrule=0.6pt,arc=2pt,left=8pt,right=8pt,top=6pt,bottom=6pt,fonttitle=\bfseries,title={#1}}

% ---------- Mã nguồn
\lstdefinelanguage{JavaScript}{
  keywords={const,let,var,function,return,if,else,for,of,in,new,throw,try,catch,async,await,class,extends,typeof,module,require,true,false,null,undefined},
  sensitive=true, comment=[l]{//}, morecomment=[s]{/*}{*/},
  morestring=[b]', morestring=[b]", morestring=[b]`}
\lstdefinelanguage{SQLx}[]{SQL}{morekeywords={REFERENCES,UNIQUE,INDEX,VIEW,CHECK,AUTOINCREMENT,IF,EXISTS,WHERE}}
\lstset{
  basicstyle=\ttfamily\footnotesize\setstretch{1.05}, backgroundcolor=\color{codebg},
  frame=single, rulecolor=\color{black!12}, framesep=6pt, xleftmargin=8pt, xrightmargin=4pt,
  keywordstyle=\color{codekw}\bfseries, commentstyle=\color{codecomment}\itshape, stringstyle=\color{codestr},
  numbers=left, numberstyle=\tiny\color{black!45}, numbersep=8pt,
  breaklines=true, breakatwhitespace=false, columns=fullflexible, keepspaces=true,
  showstringspaces=false, tabsize=2, captionpos=b, extendedchars=true,
  texcl=true}   % chú thích // được dàn trang bằng TeX: hiển thị đúng tiếng Việt
\renewcommand{\lstlistingname}{Đoạn mã}
\renewcommand{\lstlistlistingname}{Danh mục đoạn mã}

\hypersetup{pdftitle={Báo cáo Project II – Đề tài 1A},pdfauthor={Nguyễn Lê Quân – 20235199}}

% =====================================================================
\begin{document}

% ------------------------------------------------------------- TRANG BÌA
\begin{titlepage}
  \begin{tikzpicture}[remember picture,overlay]
    \draw[line width=2.4pt,color=hustred]
      ([xshift=2.2cm,yshift=-1.6cm]current page.north west) rectangle ([xshift=-1.5cm,yshift=1.6cm]current page.south east);
    \draw[line width=0.7pt,color=hustred]
      ([xshift=2.38cm,yshift=-1.78cm]current page.north west) rectangle ([xshift=-1.68cm,yshift=1.78cm]current page.south east);
  \end{tikzpicture}\par
  \centering
  {\large\bfseries ĐẠI HỌC BÁCH KHOA HÀ NỘI\par}
  \vspace{4pt}
  {\large\bfseries TRƯỜNG CÔNG NGHỆ THÔNG TIN VÀ TRUYỀN THÔNG\par}
  \vspace{4pt}
  \rule{7cm}{0.7pt}\par

  \vspace{1.3cm}
  \includegraphics[width=3.3cm]{hust_logo.png}

  \vspace{1.3cm}
  {\fontsize{24}{30}\selectfont\bfseries\color{navy} BÁO CÁO PROJECT II}\\[0.45cm]
  {\large\bfseries Đề tài 1A}\\[0.9cm]
  \begin{minipage}{\linewidth}\centering
    {\fontsize{15.5}{23}\selectfont\bfseries\color{hustred}
    PHÂN TÍCH VÀ THIẾT KẾ HỆ THỐNG HỖ TRỢ QUẢN LÝ\\ TRUNG TÂM LUYỆN THI LẤY CHỨNG CHỈ NGOẠI NGỮ\\ THEO CHUẨN QUỐC GIA, QUỐC TẾ\par}
  \end{minipage}

  \vfill
  \begin{tabular}{@{}l@{\hspace{1.2cm}}l@{}}
    \large Giảng viên hướng dẫn: & \large\bfseries TS. Nguyễn Thị Thu Hương \\[10pt]
    \large Sinh viên thực hiện:  & \large\bfseries Nguyễn Lê Quân \\[4pt]
    \large Mã số sinh viên:      & \large 20235199 \\[4pt]
    \large Lớp:                  & \large Khoa học Máy tính 07 – K68 \\[4pt]
  \end{tabular}

  \vspace{1.1cm}
  \begin{tabular}{@{}c@{}}
    \rule{5.5cm}{0.4pt}\\[-2pt]
    \small\itshape Chữ ký của giảng viên hướng dẫn
  \end{tabular}

  \vfill
  {\large\bfseries Hà Nội, tháng 11 năm 2026}
\end{titlepage}

% ------------------------------------------------------------- PHẦN ĐẦU
\pagenumbering{roman}
\setcounter{page}{2}
% ------------------------------------------------------------- LỜI CẢM ƠN
\chapter*{Lời cảm ơn}
\addcontentsline{toc}{chapter}{Lời cảm ơn}
\markboth{Lời cảm ơn}{}

Em xin gửi lời cảm ơn chân thành tới \textbf{TS. Nguyễn Thị Thu Hương} – giảng viên hướng dẫn học phần Project II – đã định hướng đề tài, góp ý về phương pháp phân tích thiết kế hướng đối tượng và cách tổ chức tài liệu trong suốt quá trình thực hiện.

Em cũng xin cảm ơn các thầy cô Trường Công nghệ Thông tin và Truyền thông, Đại học Bách khoa Hà Nội đã trang bị những kiến thức nền tảng về công nghệ phần mềm, cơ sở dữ liệu và lập trình ứng dụng web – là cơ sở để em hoàn thành đề tài này.

Do thời gian và kinh nghiệm còn hạn chế, báo cáo không tránh khỏi thiếu sót. Em rất mong nhận được ý kiến đóng góp của cô để hoàn thiện sản phẩm.

\vspace{0.6cm}
\begin{flushright}
  Hà Nội, tháng 11 năm 2026\\
  Sinh viên thực hiện\\[1.4cm]
  \textbf{Nguyễn Lê Quân}
\end{flushright}

% ------------------------------------------------------------- TÓM TẮT
\chapter*{Tóm tắt nội dung}
\addcontentsline{toc}{chapter}{Tóm tắt nội dung}
\markboth{Tóm tắt nội dung}{}

Đề tài 1A yêu cầu phân tích, thiết kế hệ thống hỗ trợ quản lý một trung tâm luyện thi chứng chỉ ngoại ngữ (IELTS, TOEIC, VSTEP, Cambridge) và lập trình bốn chức năng: \textit{quảng bá thông tin lớp học}, \textit{nhận đăng ký từ học viên}, \textit{tra cứu danh sách lớp} và \textit{chat room tư vấn}. Điểm khó của đề bài nằm ở các ràng buộc nghiệp vụ: học viên chỉ được kiểm tra trình độ sau khi nộp học phí; lớp Online mở khi có từ 10 học viên; lớp Offline mở khi đủ 20, tối đa 30 và khi vượt 30 thì việc tách lớp thuộc quyền Lãnh đạo trung tâm.

Báo cáo trình bày toàn bộ quy trình: (1) bóc tách đề bài thành 14 yêu cầu chức năng và 8 quy tắc nghiệp vụ được đánh mã; (2) mô hình hoá bằng UML với 21 biểu đồ (use case, hoạt động, tuần tự, trạng thái, lớp, ERD, kiến trúc) được viết dưới dạng mã nguồn và render cục bộ; (3) thiết kế cơ sở dữ liệu SQLite 14 bảng chuẩn 3NF, trong đó sĩ số lớp là giá trị \emph{tính toán} và các ràng buộc quan trọng được cưỡng chế ngay ở mức CSDL; (4) cài đặt ứng dụng web Node.js/Express/Socket.IO gồm cổng học viên và dashboard quản trị cho ba vai trò Tư vấn viên, Giáo vụ, Lãnh đạo; (5) kiểm thử tự động 78 ca kiểm thử (đơn vị, tích hợp API, chat realtime) đạt 100\%.

Một kết quả phân tích đáng chú ý: khi lớp Offline vượt 30 học viên, phương án “chia đôi” chỉ hợp lệ nếu tổng sĩ số đạt ít nhất 40, vì chia 31 học viên thành 16 + 15 sẽ tạo ra hai lớp không đủ điều kiện mở. Hệ thống vì vậy tự động đề xuất phương án phù hợp (chia đều, tách phần dư hoặc giữ ngoại lệ theo sức chứa phòng) để Lãnh đạo quyết định.

\noindent\textbf{Từ khoá:} phân tích thiết kế hướng đối tượng, UML, quản lý trung tâm ngoại ngữ, Node.js, Socket.IO, SQLite.


\tableofcontents
\clearpage
\listoffigures
\clearpage
\listoftables
\clearpage
\chapter*{Danh mục từ viết tắt và thuật ngữ}
\addcontentsline{toc}{chapter}{Danh mục từ viết tắt và thuật ngữ}
\markboth{Danh mục từ viết tắt}{}

\begin{xltabular}{\linewidth}{|L{3.2cm}|Y|}
  \hline
  \headrow \hd{Viết tắt / thuật ngữ} & \hd{Ý nghĩa} \\ \hline
  \endhead
  API & Application Programming Interface – giao diện lập trình ứng dụng (ở đây là REST/JSON) \\ \hline
  BR & Business Rule – quy tắc nghiệp vụ, đánh mã BR-01 … BR-08 \\ \hline
  CEFR & Common European Framework of Reference – khung tham chiếu năng lực ngôn ngữ châu Âu (A1 … C2) \\ \hline
  CSDL & Cơ sở dữ liệu \\ \hline
  ERD & Entity Relationship Diagram – biểu đồ thực thể liên kết \\ \hline
  FR / NFR & Functional / Non-functional Requirement – yêu cầu chức năng / phi chức năng \\ \hline
  GV, GVHD & Giảng viên, giảng viên hướng dẫn \\ \hline
  HV & Học viên \\ \hline
  KG & Khai giảng \\ \hline
  OOAD & Object-Oriented Analysis and Design – phân tích thiết kế hướng đối tượng \\ \hline
  Placement test & Bài kiểm tra trình độ đầu vào \\ \hline
  RBAC & Role-Based Access Control – phân quyền theo vai trò \\ \hline
  Tờ trình & Đề nghị điện tử gửi Lãnh đạo khi lớp Offline vượt 30 học viên \\ \hline
  UC & Use Case – ca sử dụng \\ \hline
  UML & Unified Modeling Language – ngôn ngữ mô hình hoá thống nhất \\ \hline
  VSTEP & Vietnamese Standardized Test of English Proficiency – bài thi đánh giá năng lực tiếng Anh theo Khung 6 bậc dùng cho Việt Nam \\ \hline
  WebSocket & Giao thức truyền thông hai chiều thời gian thực giữa trình duyệt và máy chủ \\ \hline
\end{xltabular}


% ------------------------------------------------------------- NỘI DUNG
\clearpage
\pagenumbering{arabic}
\chapter{Giới thiệu đề tài và phân tích đề bài}
\label{ch:intro}

\section{Đặt vấn đề}
Nhu cầu sở hữu chứng chỉ ngoại ngữ theo chuẩn quốc tế (IELTS, TOEIC, Cambridge) và chuẩn quốc gia (VSTEP – theo Khung năng lực ngoại ngữ 6 bậc dùng cho Việt Nam \cite{tt01}) ngày càng lớn, đặc biệt với sinh viên cần đạt chuẩn đầu ra và người đi làm cần chứng chỉ cho tuyển dụng. Các trung tâm luyện thi vì thế phải xử lý khối lượng lớn hồ sơ: quảng bá khoá học, tư vấn, thu học phí, kiểm tra đầu vào, xếp lớp theo trình độ, bố trí lịch học và phòng học.

Khi các khâu này được làm thủ công (bảng tính, tin nhắn, giấy tờ), trung tâm thường gặp các vấn đề: thông tin lớp không cập nhật kịp thời; học viên phải gọi điện nhiều lần để hỏi lịch khai giảng; giáo vụ khó nắm được lớp nào đã đủ sĩ số để mở; lớp Offline bị xếp quá chỗ ngồi; lãnh đạo không có số liệu tổng hợp để ra quyết định. Đề tài 1A đặt ra yêu cầu xây dựng hệ thống phần mềm giải quyết các vấn đề trên.

\section{Nội dung đề bài}
\begin{notebox}[Đề tài 1A – nguyên văn yêu cầu nghiệp vụ (tóm lược)]
\begin{enumerate}[leftmargin=1.6em]
  \item Trung tâm đăng tải thông tin về các lớp học, phân loại thành lớp \textbf{Online} và \textbf{Offline}.
  \item Có \textbf{phòng chat} để tư vấn trực tuyến cho người học.
  \item Khi học viên đồng ý học: hệ thống \textbf{thu học phí}, sau đó chuyển học viên sang \textbf{hệ thống kiểm tra trình độ đầu vào}; kết quả được chuyển về văn phòng trung tâm để phân loại, xếp lớp.
  \item Lớp Offline: \textbf{đủ 20 học viên} mới mở lớp, tối đa \textbf{30 học viên}; nếu vượt quá 30 thì việc tách lớp do \textbf{Lãnh đạo trung tâm} quyết định. Lớp Online có thể dưới 20 nhưng \textbf{không ít hơn 10}.
  \item Lập lịch học cho các lớp, phân phòng học cho lớp Offline, lưu trữ toàn bộ quá trình đào tạo.
  \item Phần lập trình: học một công cụ lập trình web và cài đặt 4 chức năng – \textbf{quảng bá thông tin các lớp học}, \textbf{nhận đăng ký từ học viên}, \textbf{tra cứu danh sách lớp}, \textbf{lập chat room tư vấn}.
\end{enumerate}
\end{notebox}

\section{Bóc tách đề bài}
Bước đầu tiên là chuyển từng câu chữ của đề bài thành yêu cầu có thể kiểm chứng và gán mã để truy vết xuyên suốt tài liệu, mã nguồn và kiểm thử (Bảng~\ref{tab:boc-tach}).

\begin{xltabular}{\linewidth}{|c|L{4.6cm}|Y|c|}
  \caption{Bóc tách đề bài thành yêu cầu}\label{tab:boc-tach}\\ \hline
  \headrow \hd{\#} & \hd{Câu trong đề bài} & \hd{Hiểu thành yêu cầu hệ thống} & \hd{Mã} \\ \hline
  \endfirsthead \hline
  \headrow \hd{\#} & \hd{Câu trong đề bài} & \hd{Hiểu thành yêu cầu hệ thống} & \hd{Mã} \\ \hline
  \endhead
  1 & Đăng tải thông tin lớp, phân loại Online/Offline & Danh mục khoá học (lộ trình, chuẩn đầu ra, học phí) và lớp học; mỗi lớp có thuộc tính hình thức & FR-01, 02 \\ \hline
  2 & Phòng chat tư vấn trực tuyến & Chat thời gian thực khách $\leftrightarrow$ tư vấn viên, hàng đợi, lưu lịch sử & FR-08..10 \\ \hline
  3 & Đồng ý học → thu học phí → kiểm tra đầu vào & Hồ sơ đăng ký có vòng đời; chỉ hồ sơ đã nộp phí mới được thi & FR-04..07, BR-01 \\ \hline
  4 & Kết quả kiểm tra chuyển về văn phòng để xếp lớp & Lưu kết quả, quy đổi CEFR; giáo vụ xếp lớp theo trình độ & FR-11, BR-06 \\ \hline
  5 & Offline: đủ 20 mở lớp, tối đa 30, > 30 lãnh đạo quyết định & Ngưỡng sĩ số + quy trình tờ trình phê duyệt & BR-03, 04, FR-13 \\ \hline
  6 & Online: < 20 nhưng không ít hơn 10 & Ngưỡng mở lớp online = 10 & BR-02 \\ \hline
  7 & Lập lịch, phân phòng lớp Offline & Lịch theo thứ – ca; kiểm tra trùng phòng, trùng giảng viên, sức chứa & FR-12, BR-05 \\ \hline
  8 & Lưu trữ quá trình đào tạo & Trạng thái lớp (tuyển sinh → đang học → kết thúc), danh sách học viên từng lớp & FR-12, BR-08 \\ \hline
\end{xltabular}

\subsection{Những điểm đề bài chưa nói rõ và giả định}
Trong quá trình phân tích, một số điểm đề bài không quy định cụ thể; báo cáo nêu rõ giả định để làm cơ sở thiết kế:
\begin{itemize}
  \item \textbf{Sĩ số tối đa lớp Online:} đề bài không giới hạn. Hệ thống đặt tham số cấu hình \texttt{ONLINE\_MAX = 40} (giới hạn của phòng học trực tuyến), có thể thay đổi mà không sửa mã.
  \item \textbf{Ý nghĩa của “tách lớp”:} nếu chia đôi một lớp 31 học viên sẽ được 16 + 15 – cả hai lớp đều dưới ngưỡng 20 và không được mở. Vì vậy hệ thống phân biệt hai cách tách: \emph{chia đều} (chỉ hợp lệ khi tổng $\geq 40$) và \emph{tách phần dư} sang lớp mới tiếp tục tuyển sinh (phân tích chi tiết ở mục~\ref{sec:br}).
  \item \textbf{Bài kiểm tra đầu vào:} trắc nghiệm 20 câu, 20 phút, làm một lần; kết quả quy đổi về khung CEFR \cite{cefr} để so khớp với trình độ của lớp.
  \item \textbf{Thanh toán:} cổng thanh toán (VNPay, MoMo, chuyển khoản) được mô phỏng trong prototype; giáo vụ có thể ghi nhận nộp tiền mặt tại quầy.
\end{itemize}

\section{Mục tiêu và phạm vi}
\begin{itemize}
  \item \textbf{Phân tích – thiết kế:} xây dựng đầy đủ tài liệu theo hướng đối tượng – yêu cầu, use case và đặc tả, biểu đồ hoạt động, tuần tự, trạng thái, lớp, thiết kế CSDL và kiến trúc.
  \item \textbf{Lập trình:} ứng dụng web chạy được, có CSDL thật, hiện thực 4 chức năng bắt buộc; ngoài ra cài đặt thêm luồng thu phí – kiểm tra đầu vào – xếp lớp – phê duyệt sĩ số – lập lịch phòng học để \emph{chứng minh các quy tắc nghiệp vụ hoạt động đúng}, cùng dashboard thống kê cho văn phòng.
  \item \textbf{Ngoài phạm vi:} tích hợp cổng thanh toán thật, gửi email/SMS, điểm danh – chấm điểm từng buổi học.
\end{itemize}

\section{Phương pháp và công cụ}
\begin{itemize}
  \item \textbf{Phương pháp:} phân tích thiết kế hướng đối tượng với UML 2.5 \cite{omg, fowler}; đặc tả use case theo khuôn mẫu của Cockburn \cite{cockburn}; quy trình phát triển lặp: phân tích → thiết kế → cài đặt → kiểm thử \cite{sommerville}.
  \item \textbf{Biểu đồ:} viết dưới dạng mã Mermaid \cite{mermaid} (\textit{diagram as code}) trong thư mục \texttt{docs/diagrams/src}, render cục bộ thành ảnh PNG độ phân giải cao bằng \texttt{npm run diagrams} – thay đổi thiết kế chỉ cần sửa văn bản và render lại.
  \item \textbf{Lập trình web:} Node.js 24, Express 5 \cite{express}, Socket.IO 4 \cite{socketio}, SQLite qua better-sqlite3 \cite{sqlite, bsqlite}, HTML/CSS/JavaScript thuần, Chart.js \cite{chartjs}.
  \item \textbf{Kiểm thử:} bộ kiểm thử tích hợp sẵn của Node.js (\texttt{node:test}) \cite{nodetest}.
  \item \textbf{Báo cáo:} \LaTeX{} (XeLaTeX, MiKTeX) với phông Times New Roman.
\end{itemize}

\section{Bố cục báo cáo}
Ngoài phần mở đầu, báo cáo gồm 6 chương: Chương~\ref{ch:intro} giới thiệu và bóc tách đề bài; Chương~\ref{ch:analysis} phân tích yêu cầu, quy tắc nghiệp vụ và mô hình use case; Chương~\ref{ch:design} trình bày thiết kế (hoạt động, tuần tự, trạng thái, lớp, CSDL, API); Chương~\ref{ch:impl} mô tả cài đặt và kết quả giao diện; Chương~\ref{ch:test} trình bày kiểm thử; Chương~\ref{ch:conclusion} kết luận. Phụ lục hướng dẫn cài đặt và kịch bản chạy thử.

\chapter{Phân tích yêu cầu}
\label{ch:analysis}

\section{Tác nhân của hệ thống}
\begin{xltabular}{\linewidth}{|L{3.6cm}|Y|L{5cm}|}
  \caption{Các tác nhân}\label{tab:actors}\\ \hline
  \headrow \hd{Tác nhân} & \hd{Mô tả} & \hd{Nhiệm vụ chính} \\ \hline
  \endfirsthead \hline
  \headrow \hd{Tác nhân} & \hd{Mô tả} & \hd{Nhiệm vụ chính} \\ \hline
  \endhead
  Khách / Học viên & Người quan tâm khoá học; sau khi đăng ký trở thành học viên & Xem khoá/lớp, tra cứu lớp, đăng ký, nộp học phí, làm bài kiểm tra, chat tư vấn \\ \hline
  Tư vấn viên (\texttt{CONSULTANT}) & Nhân viên tư vấn tuyển sinh & Nhận phiên chat từ hàng đợi, trả lời realtime, kết thúc phiên \\ \hline
  Giáo vụ (\texttt{ACADEMIC}) & Nhân viên văn phòng đào tạo & Thu tiền mặt, xếp lớp, lập lịch, phân phòng, mở/huỷ lớp \\ \hline
  Lãnh đạo trung tâm (\texttt{DIRECTOR}) & Ban giám đốc & Phê duyệt tờ trình vượt sĩ số; xem báo cáo; có toàn quyền của giáo vụ \\ \hline
  Cổng thanh toán & Hệ thống bên ngoài (VNPay, MoMo, ngân hàng) & Xác nhận giao dịch học phí \\ \hline
\end{xltabular}

\section{Yêu cầu chức năng}
\begin{xltabular}{\linewidth}{|l|Y|L{2.6cm}|c|}
  \caption{Danh sách yêu cầu chức năng}\label{tab:fr}\\ \hline
  \headrow \hd{Mã} & \hd{Chức năng} & \hd{Tác nhân} & \hd{Code} \\ \hline
  \endfirsthead \hline
  \headrow \hd{Mã} & \hd{Chức năng} & \hd{Tác nhân} & \hd{Code} \\ \hline
  \endhead
  FR-01 & Xem danh mục khoá học: chứng chỉ, mục tiêu đầu ra, số giờ, học phí, lộ trình, chuẩn đầu ra & Khách & CN1 \\ \hline
  FR-02 & Xem chi tiết lớp: hình thức, lịch, phòng, giảng viên, sĩ số và ngưỡng mở lớp & Khách & CN1 \\ \hline
  FR-03 & Tra cứu, lọc lớp theo từ khoá, chứng chỉ, hình thức, ca học, trình độ, ngày khai giảng, còn chỗ & Khách & CN3 \\ \hline
  FR-04 & Đăng ký học trực tuyến (thông tin cá nhân, khoá, hình thức, ca, lớp mong muốn) & Khách & CN2 \\ \hline
  FR-05 & Tra cứu hồ sơ bằng mã hồ sơ + số điện thoại & Học viên & $\checkmark$ \\ \hline
  FR-06 & Nộp học phí trực tuyến / ghi nhận tiền mặt & Học viên, Giáo vụ & $\checkmark$ \\ \hline
  FR-07 & Làm bài kiểm tra đầu vào, chấm tự động, quy đổi CEFR & Học viên & $\checkmark$ \\ \hline
  FR-08 & Mở phiên chat tư vấn, khôi phục phiên khi tải lại trang & Khách & CN4 \\ \hline
  FR-09 & Nhận phiên từ hàng đợi, trao đổi thời gian thực & Tư vấn viên & CN4 \\ \hline
  FR-10 & Kết thúc phiên, xem lịch sử hội thoại & Tư vấn viên & CN4 \\ \hline
  FR-11 & Xếp học viên vào lớp theo kết quả kiểm tra (có gợi ý) & Giáo vụ & $\checkmark$ \\ \hline
  FR-12 & Lập lịch học, phân phòng, phân giảng viên; mở, huỷ, kết thúc lớp & Giáo vụ & $\checkmark$ \\ \hline
  FR-13 & Phê duyệt tờ trình vượt sĩ số: tách lớp / giữ ngoại lệ / từ chối & Lãnh đạo & $\checkmark$ \\ \hline
  FR-14 & Dashboard: doanh thu, phễu tuyển sinh, tình trạng sĩ số, phân bố trình độ & Nhân viên & $\checkmark$ \\ \hline
\end{xltabular}
\noindent CN1–CN4 là 4 chức năng bắt buộc của đề bài; dấu $\checkmark$ là chức năng được cài đặt thêm để hoàn thiện luồng nghiệp vụ.

\section{Quy tắc nghiệp vụ}
\label{sec:br}
\begin{xltabular}{\linewidth}{|l|Y|}
  \caption{Quy tắc nghiệp vụ được chuẩn hoá}\label{tab:br}\\ \hline
  \headrow \hd{Mã} & \hd{Quy tắc} \\ \hline
  \endfirsthead \hline
  \headrow \hd{Mã} & \hd{Quy tắc} \\ \hline
  \endhead
  BR-01 & Hồ sơ chỉ được làm bài kiểm tra đầu vào khi đã thanh toán (\texttt{PAID}). Mỗi hồ sơ làm đúng một lần, thời gian 20 phút, đề được lưu lại khi phát để chống đổi đề. \\ \hline
  BR-02 & Lớp Online đủ điều kiện mở khi sĩ số $\geq 10$; sĩ số tối đa là tham số cấu hình (mặc định 40). Lớp online đầy thì từ chối xếp thêm (không lập tờ trình). \\ \hline
  BR-03 & Lớp Offline đủ điều kiện mở khi sĩ số $\geq 20$; sĩ số chuẩn $\leq 30$. \\ \hline
  BR-04 & Xếp học viên vào lớp Offline đã có 30 học viên: hệ thống \textbf{không tự xếp}, hồ sơ chuyển \texttt{WAITING\_APPROVAL}, hệ thống lập (hoặc bổ sung) \textbf{tờ trình} gửi Lãnh đạo. Mỗi lớp có tối đa một tờ trình đang chờ. \\ \hline
  BR-05 & Lớp Offline phải có phòng trước khi mở; sức chứa phòng $\geq$ sĩ số; không trùng phòng hoặc giảng viên ở cùng thứ – ca trong khoảng thời gian học giao nhau. \\ \hline
  BR-06 & Xếp lớp: lớp phải thuộc khoá đã đăng ký và nộp phí; trình độ lớp = trình độ đánh giá; ưu tiên lớp mong muốn, hình thức, ca. Xếp lệch trình độ phải ghi lý do. \\ \hline
  BR-07 & Sĩ số lớp là \textbf{giá trị tính toán} = số hồ sơ \texttt{ENROLLED} của lớp; không lưu bộ đếm để tránh sai lệch dữ liệu. \\ \hline
  BR-08 & Chỉ mở lớp (\texttt{RECRUITING} → \texttt{IN\_PROGRESS}) khi đủ ngưỡng tối thiểu, không còn tờ trình chờ, và (Offline) đã có phòng hợp lệ. Huỷ lớp thì các hồ sơ trong lớp trở về \texttt{TESTED} để xếp lớp khác. \\ \hline
\end{xltabular}

\subsection{Phân tích quy tắc tách lớp khi vượt sĩ số}
Đề bài chỉ nói “nếu số học viên vượt quá 30 thì quyết định tách lớp thuộc quyền Lãnh đạo”. Tuy nhiên phải kết hợp với BR-03 mới thấy hết bài toán. Gọi $n = 30$ là số học viên đã xếp và $k \geq 1$ là số hồ sơ đang chờ, tổng $T = n + k$. Nếu tách thành hai lớp, mỗi lớp phải có ít nhất 20 học viên mới được mở, tức là:
\[
  \left\lfloor T/2 \right\rfloor \geq 20 \iff T \geq 40 .
\]
Như vậy với $T = 31$ (trường hợp phổ biến nhất – chỉ một học viên dư) mà “chia đôi” thành $16 + 15$ thì \emph{cả hai lớp đều không đủ điều kiện mở}, đi ngược mục đích của việc tách lớp. Từ phân tích này, hệ thống đưa ra ba phương án và tự động đề xuất một phương án cho Lãnh đạo:

\begin{rulebox}[Phương án xử lý tờ trình (hàm \texttt{proposeResolution})]
\begin{enumerate}[leftmargin=1.6em]
  \item \textbf{Tách lớp – chia đều} (khi $T \geq 40$ và lớp chưa khai giảng): chia đều $T$ học viên thành $m = \max(2, \lceil T/30 \rceil)$ lớp (thường $m = 2$; ví dụ $T = 65$ thành $22 + 22 + 21$), mọi lớp đều nằm trong khoảng 20–30 nên đủ điều kiện mở; lớp mới sao chép khoá, trình độ, lịch học.
  \item \textbf{Tách lớp – tách phần dư} (khi $T < 40$ hoặc lớp đang học): giữ 30 học viên ở lớp cũ, chuyển $k$ hồ sơ sang lớp mới ở trạng thái tuyển sinh; lớp mới cần thêm $20 - k$ học viên mới mở được.
  \item \textbf{Giữ ngoại lệ}: cho lớp vượt chuẩn với $T$ học viên – \emph{chỉ khả thi khi sức chứa phòng} $\geq T$ (liên kết với BR-05).
  \item \textbf{Từ chối}: $k$ hồ sơ trở về trạng thái chờ xếp lớp để giáo vụ xếp lớp/ca khác.
\end{enumerate}
\textit{Đề xuất mặc định:} chia đều nếu $T \geq 40$; ngược lại giữ ngoại lệ nếu phòng đủ chỗ, nếu không thì tách phần dư.
\end{rulebox}

\section{Yêu cầu phi chức năng}
\begin{itemize}
  \item \textbf{Khả dụng:} giao diện tiếng Việt, responsive trên máy tính và điện thoại; có nút chuyển chế độ sáng/tối (mặc định theo hệ điều hành); thư viện giao diện được phục vụ từ chính máy chủ nên hệ thống vẫn chạy khi không có Internet (phông chữ web tự lùi về phông hệ thống).
  \item \textbf{Hiệu năng thời gian thực:} tin nhắn chat được đẩy tới bên kia qua WebSocket, độ trễ dưới 1 giây trong mạng nội bộ.
  \item \textbf{Bảo mật – riêng tư:} tra cứu hồ sơ phải khớp cả mã hồ sơ và số điện thoại; mật khẩu nhân viên băm bằng scrypt; API quản trị yêu cầu token ký HMAC có thời hạn; khoá tạm địa chỉ IP khi nhập sai mật khẩu hoặc tra sai hồ sơ quá nhiều lần; mọi truy vấn SQL đều tham số hoá; nội dung hiển thị được escape chống XSS.
  \item \textbf{Toàn vẹn dữ liệu:} khoá ngoại, ràng buộc \texttt{CHECK}, chỉ mục duy nhất có điều kiện; các thao tác nhiều bước (xếp lớp, phê duyệt, huỷ lớp) chạy trong một giao dịch.
  \item \textbf{Bảo trì:} quy tắc nghiệp vụ tập trung trong một mô-đun hàm thuần, có kiểm thử tự động bao phủ các giá trị biên.
\end{itemize}

\section{Mô hình use case}
\subsection{Biểu đồ use case tổng quát}
Hình~\ref{fig:uc0} thể hiện 10 use case chính. Bốn use case tô màu vàng là các chức năng bắt buộc phải lập trình; use case màu đỏ liên quan trực tiếp đến quy tắc sĩ số do Lãnh đạo quyết định.
\fig[0.78\linewidth]{uc_00_tong_quat}{Biểu đồ use case tổng quát}{fig:uc0}

\subsection{Biểu đồ use case phân rã}
Để tránh một biểu đồ quá lớn khiến chữ bị nhỏ, các use case được phân rã theo phân hệ (Hình~\ref{fig:uc1}–\ref{fig:uc5}).
\fig[0.85\linewidth]{uc_01_tuyen_sinh}{Phân hệ Quảng bá – Tra cứu – Đăng ký}{fig:uc1}
\fig[0.85\linewidth]{uc_02_chat}{Phân hệ Tư vấn trực tuyến}{fig:uc2}
\fig[0.85\linewidth]{uc_03_thuphi_thi}{Phân hệ Thu học phí – Kiểm tra đầu vào}{fig:uc3}
\fig[0.85\linewidth]{uc_04_xeplop}{Phân hệ Xếp lớp – Phê duyệt sĩ số}{fig:uc4}
\fig[0.85\linewidth]{uc_05_lich_phong}{Phân hệ Lịch học – Phòng học – Đào tạo}{fig:uc5}

\subsection{Đặc tả use case}
\newcommand{\ucrow}[2]{\textbf{#1} & #2 \\ \hline}
\newenvironment{ucspec}[2]{%
  \xltabular{\linewidth}{|>{\columncolor{lightgray}}L{3.3cm}|Y|}
  \caption{#1}\label{#2}\\ \hline
  \endfirsthead \hline \endhead}{\endxltabular}

\begin{ucspec}{Đặc tả UC-02 Tra cứu danh sách lớp (CN3)}{tab:uc02}
  \ucrow{Tác nhân}{Khách}
  \ucrow{Mục đích}{Tìm lớp phù hợp theo nhu cầu và xem tình trạng sĩ số}
  \ucrow{Tiền điều kiện}{Không}
  \ucrow{Luồng chính}{1. Khách nhập từ khoá và/hoặc chọn bộ lọc: chứng chỉ, hình thức, ca học, trình độ, khai giảng từ ngày, chỉ lớp còn chỗ.\newline 2. Hệ thống gửi \texttt{GET /api/classes} với các tham số.\newline 3. Hệ thống truy vấn lớp đang tuyển sinh kèm khoá học, giảng viên, phòng và sĩ số (view \texttt{v\_class\_enrollment}).\newline 4. Với mỗi lớp, hệ thống đánh giá trạng thái sĩ số theo BR-02/BR-03.\newline 5. Hiển thị danh sách thẻ lớp kèm thanh sĩ số có vạch ngưỡng mở lớp và nút “Đăng ký lớp này”.}
  \ucrow{Luồng thay thế}{3a. Không có kết quả: hiển thị thông báo và gợi ý nới bộ lọc.\newline 1a. Khách đổi bộ lọc: kết quả cập nhật tức thì (debounce 220 ms).}
  \ucrow{Hậu điều kiện}{Không thay đổi dữ liệu}
\end{ucspec}

\begin{ucspec}{Đặc tả UC-03 Đăng ký học (CN2)}{tab:uc03}
  \ucrow{Tác nhân}{Khách}
  \ucrow{Tiền điều kiện}{Khoá học đang được công bố}
  \ucrow{Luồng chính}{1. Khách mở form (trực tiếp hoặc từ nút “Đăng ký lớp này”, form tự điền khoá, hình thức, ca, lớp).\newline 2. Nhập họ tên, ngày sinh, SĐT, email, khoá học, hình thức, ca mong muốn, lớp mong muốn, ghi chú.\newline 3. Hệ thống kiểm tra hợp lệ phía trình duyệt rồi gửi \texttt{POST /api/registrations}.\newline 4. Máy chủ kiểm tra lại toàn bộ dữ liệu; tìm học viên cũ theo email/SĐT, nếu chưa có thì tạo mới.\newline 5. Tạo hồ sơ trạng thái \texttt{PENDING\_PAYMENT}, cấp mã \texttt{DK-YYYY-NNNN}.\newline 6. Hiển thị mã hồ sơ, học phí và nút “Nộp học phí ngay”.}
  \ucrow{Ngoại lệ}{3a/4a. Dữ liệu sai: trả 400 kèm lỗi theo từng trường, form tô đỏ trường lỗi.\newline 4b. Học viên đã có hồ sơ hiệu lực cho cùng khoá: trả 409 kèm mã hồ sơ cũ.\newline 4c. Email và SĐT thuộc hai học viên khác nhau: trả 409.\newline 4d. Lớp mong muốn không thuộc khoá hoặc khác hình thức: trả 400.}
  \ucrow{Hậu điều kiện}{Có bản ghi \texttt{students} (nếu mới) và \texttt{registrations}}
\end{ucspec}

\begin{ucspec}{Đặc tả UC-05 Kiểm tra trình độ đầu vào}{tab:uc05}
  \ucrow{Tác nhân}{Học viên}
  \ucrow{Tiền điều kiện}{Hồ sơ ở trạng thái \texttt{PAID} (BR-01)}
  \ucrow{Luồng chính}{1. Học viên tra cứu hồ sơ (mã + SĐT) và bấm “Bắt đầu làm bài”.\newline 2. Hệ thống tạo lượt thi, chọn ngẫu nhiên 20 câu cân bằng 5 mức độ khó, lưu danh sách câu đã phát; trả câu hỏi \emph{không kèm đáp án}.\newline 3. Học viên làm bài, đồng hồ đếm ngược 20 phút; câu trả lời được lưu tạm ở trình duyệt.\newline 4. Học viên nộp bài; hệ thống chấm theo đề đã lưu, quy đổi CEFR, chuyển hồ sơ sang \texttt{TESTED}.\newline 5. Kết quả hiển thị cho học viên và xuất hiện ở văn phòng trong danh sách “Chờ xếp lớp”.}
  \ucrow{Ngoại lệ}{1a. Chưa nộp học phí: 403 \texttt{PAYMENT\_REQUIRED}.\newline 4a. Đã nộp trước đó: 409.\newline 4b. Quá giờ (cho phép trễ 60 giây do mạng): 410.\newline 3a. Tải lại trang giữa chừng: tiếp tục đúng đề cũ, thời gian vẫn tính.}
\end{ucspec}

\begin{ucspec}{Đặc tả UC-06 Tư vấn trực tuyến (CN4)}{tab:uc06}
  \ucrow{Tác nhân}{Khách (chính), Tư vấn viên}
  \ucrow{Luồng chính}{1. Khách mở widget chat, nhập tên, liên hệ, chủ đề.\newline 2. Hệ thống tạo phiên \texttt{WAITING}, gửi lời chào tự động, đẩy sự kiện cập nhật hàng đợi tới các tư vấn viên đang trực tuyến.\newline 3. Tư vấn viên bấm “Nhận phiên”; hệ thống cập nhật có điều kiện (chỉ một người nhận được), phiên chuyển \texttt{ACTIVE}, gửi tin hệ thống “Tư vấn viên X đã tham gia”.\newline 4. Hai bên trao đổi; mỗi tin nhắn được lưu CSDL rồi phát tới phòng chat; có chỉ báo “đang soạn tin”.\newline 5. Tư vấn viên (hoặc khách) kết thúc; phiên chuyển \texttt{CLOSED}.}
  \ucrow{Luồng thay thế}{1a. Khách tải lại trang: widget dùng \texttt{sessionId} đã lưu để khôi phục phiên và toàn bộ lịch sử.\newline 3a. Tư vấn viên khác đã nhận: thông báo lỗi “phiên đã được nhận”.}
  \ucrow{Ràng buộc}{Tư vấn viên phải nhận phiên mới được trả lời; socket lạ không gửi tin vào phiên của người khác được.}
\end{ucspec}

\begin{ucspec}{Đặc tả UC-07 Xếp lớp và UC-08 Phê duyệt tờ trình}{tab:uc07}
  \ucrow{Tác nhân}{Giáo vụ (UC-07), Lãnh đạo trung tâm (UC-08)}
  \ucrow{Tiền điều kiện}{Hồ sơ \texttt{TESTED} (UC-07); tờ trình \texttt{PENDING} (UC-08)}
  \ucrow{Luồng chính UC-07}{1. Giáo vụ mở hồ sơ, hệ thống gợi ý danh sách lớp đã xếp hạng theo BR-06.\newline 2. Giáo vụ chọn lớp. 3. Hệ thống áp dụng \texttt{placementDecision}: Online còn chỗ hoặc Offline < 30 → \texttt{ENROLLED}.}
  \ucrow{Rẽ nhánh UC-07}{3a. Offline đã đủ 30 → hồ sơ \texttt{WAITING\_APPROVAL}, tạo/bổ sung tờ trình (BR-04).\newline 3b. Online đầy, sai khoá, lớp đã đóng → từ chối kèm lý do.\newline 1a. Lệch trình độ → bắt buộc nhập lý do, lưu vào hồ sơ.}
  \ucrow{Luồng chính UC-08}{1. Lãnh đạo xem tờ trình: sĩ số hiện tại, số hồ sơ chờ, sức chứa phòng, danh sách hồ sơ chờ và phương án đề xuất.\newline 2. Chọn quyết định (tách lớp / giữ ngoại lệ / từ chối) và ghi chú.\newline 3. Hệ thống thực thi trong \textbf{một giao dịch}; ghi người duyệt, thời điểm, lớp mới (nếu tách).}
  \ucrow{Ngoại lệ UC-08}{Giữ ngoại lệ khi phòng không đủ chỗ → 422; tờ trình đã xử lý → 409; người không phải Lãnh đạo → 403.}
\end{ucspec}

\begin{ucspec}{Đặc tả UC-04 Thanh toán học phí}{tab:uc04}
  \ucrow{Tác nhân}{Học viên, Cổng thanh toán; hoặc Giáo vụ (thu tiền mặt)}
  \ucrow{Tiền điều kiện}{Hồ sơ ở trạng thái \texttt{PENDING\_PAYMENT}}
  \ucrow{Luồng chính}{1. Học viên tra cứu hồ sơ bằng mã hồ sơ + SĐT.\newline 2. Chọn phương thức (VNPay, MoMo, chuyển khoản) và xác nhận.\newline 3. Cổng thanh toán xác nhận giao dịch (mô phỏng).\newline 4. Trong một giao dịch CSDL: tạo bản ghi \texttt{payments} với số tiền bằng học phí khoá, cập nhật hồ sơ sang \texttt{PAID}.}
  \ucrow{Luồng thay thế}{2a. Học viên nộp tiền mặt tại quầy: giáo vụ bấm “Thu tiền mặt”, hệ thống ghi nhận người thu.}
  \ucrow{Ngoại lệ}{Hồ sơ đã thanh toán → 409; hồ sơ đã huỷ → 409; tra sai mã/SĐT nhiều lần → 429 (khoá tạm).}
\end{ucspec}

\begin{ucspec}{Đặc tả UC-09 Lập lịch, phân phòng, mở lớp}{tab:uc09}
  \ucrow{Tác nhân}{Giáo vụ}
  \ucrow{Tiền điều kiện}{Lớp ở trạng thái \texttt{RECRUITING} hoặc \texttt{IN\_PROGRESS}}
  \ucrow{Luồng chính}{1. Giáo vụ mở chi tiết lớp; hệ thống hiển thị sĩ số, danh sách học viên và các lý do lớp chưa đủ điều kiện mở (nếu có).\newline 2. Chọn giảng viên, phòng học (lớp Offline) và các buổi học theo thứ – ca.\newline 3. Hệ thống kiểm tra BR-05: sức chứa phòng, trùng phòng, trùng giảng viên trong khoảng thời gian học giao nhau; hợp lệ thì lưu.\newline 4. Giáo vụ bấm “Mở lớp”; hệ thống kiểm tra BR-08 rồi chuyển lớp sang \texttt{IN\_PROGRESS}.}
  \ucrow{Ngoại lệ}{3a. Trùng lịch → 422 kèm tên lớp và buổi bị trùng.\newline 4a. Chưa đủ sĩ số tối thiểu, còn tờ trình chờ, hoặc lớp Offline chưa có phòng → 422 kèm danh sách lý do.}
  \ucrow{Luồng thay thế}{Huỷ lớp đang tuyển sinh: các hồ sơ trong lớp trở về \texttt{TESTED}, tờ trình đang chờ (nếu có) được đóng.}
\end{ucspec}

\section{Ma trận phân quyền}
\begin{table}[H]\centering
\caption{Ma trận phân quyền theo vai trò}\label{tab:rbac}
\begin{tabular}{|l|c|c|c|c|}
  \hline
  \headrow \hd{Chức năng} & \hd{Khách} & \hd{Tư vấn viên} & \hd{Giáo vụ} & \hd{Lãnh đạo} \\ \hline
  Xem, tra cứu, đăng ký, nộp phí, làm bài & $\checkmark$ & $\checkmark$ & $\checkmark$ & $\checkmark$ \\ \hline
  Trả lời chat tư vấn & & $\checkmark$ & & $\checkmark$ \\ \hline
  Xem dashboard, danh sách hồ sơ, lớp & & $\checkmark$ & $\checkmark$ & $\checkmark$ \\ \hline
  Thu tiền mặt, xếp lớp, lập lịch, mở/huỷ lớp & & & $\checkmark$ & $\checkmark$ \\ \hline
  Phê duyệt tờ trình vượt sĩ số & & & & $\checkmark$ \\ \hline
\end{tabular}
\end{table}

\chapter{Thiết kế hệ thống}
\label{ch:design}

\section{Kiến trúc tổng thể}
Hệ thống theo mô hình client – server ba tầng (Hình~\ref{fig:arch}). Trình duyệt chạy hai ứng dụng: \textit{Cổng học viên} và \textit{Dashboard quản trị}. Máy chủ Node.js cung cấp REST API (Express) và kênh WebSocket (Socket.IO) cho chat; toàn bộ dữ liệu lưu trong một tệp SQLite.
\fig[0.95\linewidth]{arch_01}{Kiến trúc tổng thể và luồng giao tiếp}{fig:arch}

Phía máy chủ được tổ chức phân tầng (Hình~\ref{fig:layers}): tầng \textbf{routes} chỉ nhận yêu cầu và trả kết quả; tầng \textbf{services} điều phối truy cập CSDL và giao dịch; tầng \textbf{domain} (\texttt{rules.js}) chứa các quy tắc nghiệp vụ dưới dạng \emph{hàm thuần} – không đọc/ghi CSDL – nên có thể kiểm thử độc lập và dùng chung cho mọi luồng (API, socket, dữ liệu mẫu).
\fig[0.8\linewidth]{cls_02_layers}{Biểu đồ lớp thiết kế phân tầng phía máy chủ}{fig:layers}

\section{Biểu đồ hoạt động}
Các biểu đồ hoạt động dùng làn (swimlane) để phân định trách nhiệm giữa tác nhân và hệ thống.
\fig[0.92\linewidth]{act_01_dangky_thuphi_thi}{Hoạt động: Đăng ký – Thu học phí – Kiểm tra đầu vào}{fig:act1}
Hình~\ref{fig:act1} cho thấy điểm kiểm soát BR-01: nhánh “Hồ sơ đã PAID?” nằm trong làn hệ thống, nên kể cả khi người dùng gọi API trực tiếp (bỏ qua giao diện) cũng không thể làm bài khi chưa nộp phí.
\fig[0.92\linewidth]{act_02_xeplop_pheduyet}{Hoạt động: Xếp lớp và xử lý vượt sĩ số}{fig:act2}
Trong Hình~\ref{fig:act2}, lớp Online và Offline đi theo hai nhánh khác nhau: lớp Online đầy thì từ chối và yêu cầu chọn lớp khác, còn lớp Offline đủ 30 thì chuyển sang luồng tờ trình cho Lãnh đạo – đúng tinh thần đề bài “quyết định tách lớp thuộc quyền Lãnh đạo”.
\fig[0.85\linewidth]{act_03_chat}{Hoạt động: Tư vấn trực tuyến}{fig:act3}
\fig[0.95\linewidth]{act_04_tracuu}{Hoạt động: Tra cứu danh sách lớp}{fig:act4}

\section{Biểu đồ tuần tự}
\fig[0.95\linewidth]{seq_01_dangky}{Tuần tự: Đăng ký, thanh toán và làm bài kiểm tra}{fig:seq1}
\fig[0.95\linewidth]{seq_02_tracuu}{Tuần tự: Tra cứu danh sách lớp}{fig:seq2}
\fig[0.95\linewidth]{seq_03_chat}{Tuần tự: Chat tư vấn thời gian thực qua Socket.IO}{fig:seq3}
\fig[0.95\linewidth]{seq_04_xeplop_pheduyet}{Tuần tự: Xếp lớp và phê duyệt tờ trình}{fig:seq4}

\section{Biểu đồ trạng thái}
Hệ thống tách bạch hai loại trạng thái của lớp học:
\begin{itemize}
  \item \textbf{Trạng thái vòng đời} được lưu trong CSDL: \texttt{RECRUITING}, \texttt{IN\_PROGRESS}, \texttt{COMPLETED}, \texttt{CANCELLED}. Chỉ thay đổi bằng thao tác có chủ đích (mở lớp, kết thúc, huỷ).
  \item \textbf{Trạng thái sĩ số} được \emph{tính toán} mỗi lần đọc từ sĩ số và số hồ sơ chờ: \texttt{BELOW\_MIN}, \texttt{ELIGIBLE}, \texttt{FULL}, \texttt{PENDING\_APPROVAL}, \texttt{OVER\_CAPACITY} (các trạng thái con trong Hình~\ref{fig:st1}). Nhờ vậy không bao giờ xảy ra tình huống “lớp ghi là đủ điều kiện mở nhưng thực tế đã có học viên rút”.
\end{itemize}
\fig[0.98\linewidth]{st_01_lop}{Biểu đồ trạng thái của Lớp học}{fig:st1}
\fig[0.98\linewidth]{st_02_hoso}{Biểu đồ trạng thái của Hồ sơ đăng ký}{fig:st2}
\fig[0.8\linewidth]{st_03_chat}{Biểu đồ trạng thái của Phiên tư vấn}{fig:st3}
Ràng buộc giữa trạng thái hồ sơ và lớp được cưỡng chế ở mức CSDL: hồ sơ có \texttt{class\_id} khác rỗng \emph{khi và chỉ khi} ở trạng thái \texttt{ENROLLED} hoặc \texttt{WAITING\_APPROVAL}.

\section{Biểu đồ lớp miền nghiệp vụ}
\fig[0.98\linewidth]{cls_01_domain}{Biểu đồ lớp miền nghiệp vụ}{fig:domain}
Lớp \texttt{Class} có các phương thức \texttt{minSize()}, \texttt{maxSize()}, \texttt{capacityState()}, \texttt{canOpen()} – trong cài đặt, các phương thức này được hiện thực bởi các hàm thuần trong \texttt{rules.js}. Quan hệ tự thân “tách từ” (\texttt{parent\_class\_id}) cho phép truy vết lớp mới sinh ra từ quyết định tách lớp của Lãnh đạo.

\section{Thiết kế cơ sở dữ liệu}
\subsection{Lược đồ thực thể – liên kết}
\fig[0.98\linewidth]{erd_01}{Lược đồ ERD – 14 bảng}{fig:erd}

\subsection{Từ điển dữ liệu các bảng chính}
\begin{xltabular}{\linewidth}{|L{3.4cm}|L{3.2cm}|Y|}
  \caption{Từ điển dữ liệu (rút gọn)}\label{tab:dict}\\ \hline
  \headrow \hd{Bảng} & \hd{Cột chính} & \hd{Ràng buộc / ý nghĩa} \\ \hline
  \endfirsthead \hline
  \headrow \hd{Bảng} & \hd{Cột chính} & \hd{Ràng buộc / ý nghĩa} \\ \hline
  \endhead
  \texttt{courses} & code, certificate, target, tuition\_fee, syllabus & \texttt{code} duy nhất; \texttt{certificate} $\in$ \{IELTS, TOEIC, VSTEP, CAMBRIDGE\}; lộ trình lưu JSON \\ \hline
  \texttt{classes} & code, course\_id, level, form, start/end\_date, teacher\_id, room\_id, status, parent\_class\_id & \texttt{form} $\in$ \{ONLINE, OFFLINE\}; \texttt{CHECK(end\_date > start\_date)}; \texttt{CHECK(form = 'OFFLINE' OR room\_id IS NULL)} – lớp online không có phòng \\ \hline
  \texttt{class\_schedules} & class\_id, weekday (2–8), shift & \texttt{UNIQUE(class\_id, weekday, shift)}; ca $\in$ \{MORNING, AFTERNOON, EVENING\} \\ \hline
  \texttt{rooms}, \texttt{teachers} & capacity, campus / title, specialty & \texttt{capacity > 0} \\ \hline
  \texttt{students} & full\_name, dob, phone, email & \texttt{phone}, \texttt{email} duy nhất – tái sử dụng học viên cũ \\ \hline
  \texttt{registrations} & reg\_code, student\_id, course\_id, desired\_*, class\_id, status, placement\_note & 6 trạng thái; \texttt{CHECK((status IN ('ENROLLED','WAITING\_APPROVAL')) = (class\_id IS NOT NULL))}; chỉ mục duy nhất có điều kiện: một hồ sơ hiệu lực / học viên / khoá \\ \hline
  \texttt{payments} & registration\_id, amount, method, txn\_code, collected\_by & \texttt{registration\_id} duy nhất (mỗi hồ sơ thanh toán 1 lần); \texttt{txn\_code} duy nhất \\ \hline
  \texttt{placement\_questions} & skill, difficulty, options, answer\_index & 5 mức độ khó; đáp án không bao giờ gửi ra trình duyệt \\ \hline
  \texttt{placement\_attempts} & registration\_id, question\_ids, score, level, started/submitted\_at & \texttt{registration\_id} duy nhất (thi 1 lần); lưu đề đã phát \\ \hline
  \texttt{approval\_requests} & class\_id, status, proposed\_action, decision, decided\_by, result\_class\_id & Chỉ mục duy nhất có điều kiện \texttt{WHERE status = 'PENDING'}: một tờ trình chờ / lớp \\ \hline
  \texttt{staff\_users} & username, password\_hash, role & role $\in$ \{CONSULTANT, ACADEMIC, DIRECTOR\}; mật khẩu băm scrypt + salt \\ \hline
  \texttt{chat\_sessions}, \texttt{chat\_messages} & session\_id (UUID), status, consultant\_id / sender\_type, body & \texttt{body} 1–2000 ký tự; xoá phiên thì xoá tin nhắn (\texttt{ON DELETE CASCADE}) \\ \hline
\end{xltabular}

\subsection{Sĩ số là giá trị tính toán}
Thay vì lưu cột \texttt{current\_enrolled} trong bảng \texttt{classes} (dễ lệch khi có lỗi giữa chừng), sĩ số được tính qua view:
\begin{lstlisting}[language=SQLx,caption={View tính sĩ số và hồ sơ chờ duyệt của từng lớp},label=lst:view]
CREATE VIEW v_class_enrollment AS
SELECT c.class_id,
  (SELECT COUNT(*) FROM registrations r
    WHERE r.class_id = c.class_id AND r.status = 'ENROLLED')         AS enrolled,
  (SELECT COUNT(*) FROM registrations r
    WHERE r.class_id = c.class_id AND r.status = 'WAITING_APPROVAL') AS waiting
FROM classes c;
\end{lstlisting}
Kết hợp chỉ mục \texttt{ix\_reg\_class\_status(class\_id, status)}, truy vấn này vẫn nhanh với quy mô một trung tâm (hàng nghìn hồ sơ).

\section{Thiết kế API}
\begin{xltabular}{\linewidth}{|l|L{6.3cm}|Y|}
  \caption{Các API chính}\label{tab:api}\\ \hline
  \headrow \hd{Phương thức} & \hd{Đường dẫn} & \hd{Chức năng – quyền} \\ \hline
  \endfirsthead \hline
  \headrow \hd{Phương thức} & \hd{Đường dẫn} & \hd{Chức năng – quyền} \\ \hline
  \endhead
  GET & \texttt{/api/courses}, \texttt{/api/courses/:id} & CN1 – danh mục, chi tiết khoá + lớp (công khai) \\ \hline
  GET & \texttt{/api/classes?q\&certificate\&form\&shift\&level\&from\&available} & CN3 – tra cứu lớp (công khai) \\ \hline
  POST & \texttt{/api/registrations} & CN2 – đăng ký (công khai) \\ \hline
  POST & \texttt{/api/registrations/lookup} & Tra cứu hồ sơ bằng mã + SĐT \\ \hline
  POST & \texttt{/api/registrations/:code/pay} & Thanh toán trực tuyến (mô phỏng) \\ \hline
  POST & \texttt{/api/placement/:code/start}, \texttt{/submit} & Làm bài / nộp bài kiểm tra (BR-01) \\ \hline
  POST & \texttt{/api/auth/login} & Đăng nhập nhân viên, trả token \\ \hline
  GET & \texttt{/api/admin/dashboard} & Số liệu tổng hợp – mọi nhân viên \\ \hline
  GET & \texttt{/api/admin/registrations/:id/suggestions} & Gợi ý lớp – Giáo vụ, Lãnh đạo \\ \hline
  POST & \texttt{/api/admin/registrations/:id/enroll} & Xếp lớp (BR-04, BR-06) – Giáo vụ, Lãnh đạo \\ \hline
  POST & \texttt{/api/admin/registrations/:id/cash-payment} & Thu tiền mặt – Giáo vụ, Lãnh đạo \\ \hline
  POST/PUT & \texttt{/api/admin/classes}, \texttt{/:id/arrangement} & Tạo lớp, lập lịch – phân phòng (BR-05) \\ \hline
  POST & \texttt{/api/admin/classes/:id/open|cancel|complete} & Mở / huỷ / kết thúc lớp (BR-08) \\ \hline
  GET & \texttt{/api/admin/approvals} & Danh sách tờ trình kèm phương án đề xuất \\ \hline
  POST & \texttt{/api/admin/approvals/:id/decide} & Ra quyết định – \textbf{chỉ Lãnh đạo} \\ \hline
  GET & \texttt{/api/admin/chat/sessions[/:id/messages]} & Lịch sử chat – nhân viên \\ \hline
\end{xltabular}
Mã lỗi HTTP được dùng nhất quán: 400 dữ liệu sai, 401 chưa đăng nhập, 403 không đủ quyền hoặc vi phạm BR-01, 404 không tìm thấy, 409 xung đột trạng thái (trùng hồ sơ, đã thanh toán, tờ trình đã xử lý), 410 hết giờ làm bài, 422 vi phạm quy tắc nghiệp vụ (lệch trình độ, trùng lịch, chưa đủ điều kiện mở lớp).

\section{Thiết kế giao diện}
\begin{itemize}
  \item \textbf{Cổng học viên} là trang đơn gồm các khu vực: giới thiệu và quy định mở lớp, danh mục khoá học (CN1), tra cứu lớp (CN3), form đăng ký (CN2), cửa sổ tra cứu hồ sơ – thanh toán – làm bài, và widget chat nổi ở góc màn hình (CN4).
  \item \textbf{Dashboard quản trị} có thanh điều hướng bên trái với 5 mục: Tổng quan, Hồ sơ \& xếp lớp, Lớp học \& lịch, Phê duyệt sĩ số, Tư vấn trực tuyến; huy hiệu đỏ báo số việc đang chờ (cập nhật realtime).
  \item Màu sắc nhất quán cho trạng thái sĩ số: vàng – chưa đủ ngưỡng, xanh lá – đủ điều kiện mở, xanh dương – đủ 30, đỏ – chờ Lãnh đạo duyệt. Thanh sĩ số có vạch đánh dấu ngưỡng mở lớp (10 hoặc 20) giúp nhìn là biết lớp còn thiếu bao nhiêu học viên.
\end{itemize}

\chapter{Cài đặt hệ thống}
\label{ch:impl}

\section{Công nghệ sử dụng}
\begin{table}[H]\centering
\caption{Công nghệ và lý do lựa chọn}\label{tab:tech}
\begin{tabularx}{\linewidth}{|L{3.3cm}|L{3.4cm}|Y|}
  \hline
  \headrow \hd{Thành phần} & \hd{Công nghệ} & \hd{Lý do lựa chọn} \\ \hline
  Máy chủ & Node.js 24, Express 5 & Một ngôn ngữ (JavaScript) cho cả client và server; Express gọn nhẹ, tự chuyển lỗi đồng bộ tới bộ xử lý lỗi tập trung \\ \hline
  Thời gian thực & Socket.IO 4 & Hỗ trợ phòng (room), xác nhận (ack), tự kết nối lại – phù hợp chat có hàng đợi \\ \hline
  Cơ sở dữ liệu & SQLite 3 + better-sqlite3 & Không cần cài máy chủ CSDL, một tệp duy nhất; API đồng bộ và \texttt{transaction()} giúp viết giao dịch rõ ràng; hỗ trợ partial index, view, CHECK \\ \hline
  Giao diện & HTML5, CSS3, JavaScript thuần & Không cần bước build; hệ thống thiết kế tự viết (biến CSS, chế độ sáng/tối) \\ \hline
  Biểu đồ thống kê & Chart.js 4 & Phục vụ cục bộ từ \texttt{node\_modules}, chạy offline \\ \hline
  Kiểm thử & \texttt{node:test}, socket.io-client & Có sẵn trong Node.js, không cần thư viện ngoài \\ \hline
  Biểu đồ UML & Mermaid CLI & Biểu đồ dạng mã, render cục bộ bằng trình duyệt có sẵn \\ \hline
\end{tabularx}
\end{table}

\section{Cấu trúc mã nguồn}
\begin{figure}[H]
\begin{PlainCode}
ProjectII/
|-- server/
|   |-- index.js            khởi động HTTP + Socket.IO, nạp dữ liệu mẫu khi CSDL trống
|   |-- app.js              cấu hình Express, định tuyến, xử lý lỗi tập trung
|   |-- config.js           ngưỡng sĩ số (10 / 20 / 30), ca học, thời gian thi
|   |-- domain/rules.js     BUSINESS RULES: hàm thuần BR-01..BR-08
|   |-- services/           catalog, registration, placement, enrollment, chat, stats, auth
|   |-- routes/             public.js (cổng học viên), admin.js (văn phòng)
|   |-- sockets/chat.js     sự kiện chat realtime
|   |-- middleware/auth.js  kiểm tra token + vai trò
|   `-- db/                 schema.sql, connection.js, seed.js, questions.js
|-- public/                 index.html (cổng học viên), admin.html (dashboard), css/, js/
|-- tests/                  rules.test.js, api.test.js, chat.test.js
|-- docs/                   analysis/*.md, diagrams/src/*.mmd -> diagrams/png/*.png
|-- tools/                  render-diagrams.js, screenshots.js
`-- report/                 báo cáo LaTeX (main.tex, chapters/, images/)
\end{PlainCode}
\captionof{lstlisting}{Cấu trúc thư mục dự án}\label{lst:tree}
\end{figure}
Toàn bộ dự án gồm khoảng 4\,200 dòng mã (JavaScript, SQL, HTML, CSS), không tính thư viện.

\section{Cài đặt bộ quy tắc nghiệp vụ}
Mọi quy tắc sĩ số nằm trong \texttt{server/domain/rules.js}. Đoạn mã~\ref{lst:decision} là hàm quyết định khi giáo vụ xếp một hồ sơ vào lớp – ánh xạ trực tiếp BR-02, BR-04, BR-06.
\begin{lstlisting}[language=JavaScript,caption={Hàm \texttt{placementDecision} (rút gọn)},label=lst:decision]
function placementDecision({ cls, enrolled, registration, studentLevel, reason }) {
  if (registration.status !== 'TESTED')
    return { action: 'REJECT', code: 'INVALID_STATUS' };
  if (cls.course_id !== registration.course_id)
    return { action: 'REJECT', code: 'WRONG_COURSE' };   // phải đúng khoá đã nộp phí
  if (!['RECRUITING', 'IN_PROGRESS'].includes(cls.status))
    return { action: 'REJECT', code: 'CLASS_CLOSED' };
  if (studentLevel && cls.level !== studentLevel && !(reason && reason.trim().length >= 5))
    return { action: 'REJECT', code: 'LEVEL_MISMATCH' };  // BR-06
  const { max } = limitsOf(cls.form);                     // ONLINE 40 / OFFLINE 30
  if (cls.form === 'ONLINE')
    return enrolled >= max ? { action: 'REJECT', code: 'ONLINE_FULL' }   // BR-02
                           : { action: 'ENROLL', code: 'OK' };
  if (enrolled >= max)
    return { action: 'NEEDS_APPROVAL', code: 'OVER_30' };  // BR-04: không tự xếp
  return { action: 'ENROLL', code: 'OK' };
}
\end{lstlisting}

Đoạn mã~\ref{lst:propose} hiện thực phân tích ở mục~\ref{sec:br}: chỉ chia đều khi tổng sĩ số đủ để cả hai lớp đạt ngưỡng 20.
\begin{lstlisting}[language=JavaScript,caption={Hàm \texttt{proposeResolution} – phương án xử lý tờ trình},label=lst:propose]
function proposeResolution({ enrolled, waiting, roomCapacity = null, inProgress = false }) {
  const { min } = CAPACITY.OFFLINE;                 // 20
  const total = enrolled + waiting;
  const even = !inProgress && total >= 2 * min;     // 40: hai lớp đều >= 20
  const split = even
    ? { mode: 'EVEN', sizes: [Math.ceil(total / 2), Math.floor(total / 2)] }
    : { mode: 'REMAINDER', sizes: [enrolled, waiting] };   // lớp mới tiếp tục tuyển sinh
  const keepFeasible = roomCapacity == null || roomCapacity >= total;  // BR-05
  const recommended = even || !keepFeasible ? 'SPLIT' : 'KEEP';
  return { total, recommended, options: { SPLIT: split, KEEP: { feasible: keepFeasible }, REJECT: {} } };
}
\end{lstlisting}

\section{Cài đặt giao dịch xếp lớp và phê duyệt}
Các thao tác thay đổi nhiều bảng được bọc trong \texttt{db.tx()} (giao dịch của better-sqlite3): nếu bất kỳ bước nào ném lỗi, toàn bộ thay đổi bị huỷ. Khi Lãnh đạo chọn \emph{tách lớp}, hệ thống trong cùng một giao dịch: tạo lớp mới (sao chép khoá, trình độ, hình thức, lịch học, gắn \texttt{parent\_class\_id}); tự tìm giảng viên và phòng còn trống không trùng lịch (BR-05); chuyển các hồ sơ chờ sang lớp mới; nếu chia đều thì chuyển thêm những học viên được xếp muộn nhất; cuối cùng cập nhật tờ trình.
\begin{lstlisting}[language=JavaScript,caption={Xếp lớp: lập hoặc bổ sung tờ trình trong cùng giao dịch (rút gọn)},label=lst:enroll]
function enroll(regId, classId, reason, staffId) {
  return db.tx(() => {
    const reg = d.prepare('SELECT * FROM registrations WHERE registration_id = ?').get(regId);
    const cls = rawClass(classId);
    const { enrolled, waiting } = counts(classId);           // từ view v\_class\_enrollment
    const decision = rules.placementDecision({ cls, enrolled, registration: reg,
                                               studentLevel: assessedLevel(regId), reason });
    if (decision.action === 'REJECT') throw unprocessable(decision.code, decision.message);
    if (decision.action === 'ENROLL') { /* UPDATE ... status = 'ENROLLED' */ return ...; }
    // NEEDS\_APPROVAL: lập hoặc bổ sung tờ trình
    d.prepare(`UPDATE registrations SET class_id = ?, status = 'WAITING_APPROVAL' ...`).run(...);
    const pending = d.prepare("SELECT request_id FROM approval_requests " +
                              "WHERE class_id = ? AND status = 'PENDING'").get(classId);
    if (pending) { ... }      // bổ sung vào tờ trình đang chờ
    else { ... }             // INSERT approval\_requests (PENDING)
  });
}
\end{lstlisting}
Chỉ mục duy nhất có điều kiện \texttt{ux\_pending\_request ON approval\_requests(class\_id) WHERE status = 'PENDING'} bảo đảm kể cả khi hai giáo vụ thao tác đồng thời cũng không thể sinh hai tờ trình cho cùng một lớp.

\section{Cài đặt chat thời gian thực}
\begin{table}[H]\centering
\caption{Các sự kiện Socket.IO}\label{tab:events}
\begin{tabularx}{\linewidth}{|l|L{2.6cm}|Y|}
  \hline
  \headrow \hd{Sự kiện} & \hd{Phía gửi} & \hd{Xử lý} \\ \hline
  \texttt{chat:start} & Khách & Tạo phiên WAITING, gửi lời chào, phát \texttt{queue:update} tới phòng \texttt{consultants} \\ \hline
  \texttt{chat:resume} & Khách / NV & Tham gia lại phòng \texttt{session:<id>}, trả lịch sử tin nhắn \\ \hline
  \texttt{chat:claim} & Tư vấn viên & \texttt{UPDATE ... WHERE status = 'WAITING'} – chỉ một người nhận được; phiên ACTIVE \\ \hline
  \texttt{chat:message} & Hai phía & Kiểm tra quyền (khách đúng phiên, tư vấn viên đã nhận phiên), lưu CSDL rồi phát tới phòng \\ \hline
  \texttt{chat:typing} & Hai phía & Chuyển tiếp chỉ báo “đang soạn tin” \\ \hline
  \texttt{chat:close} & Hai phía & Phiên CLOSED, gửi tin hệ thống, phát \texttt{chat:closed} \\ \hline
\end{tabularx}
\end{table}
Tư vấn viên kết nối socket kèm token trong \texttt{handshake.auth}; máy chủ xác thực token và đưa socket vào phòng \texttt{consultants}. Mọi handler được bọc để lỗi nghiệp vụ trả về qua \texttt{ack} thay vì làm đứt kết nối.

\section{Bảo mật}
\begin{itemize}
  \item \textbf{Mật khẩu:} \texttt{scrypt} với salt ngẫu nhiên 16 byte; so sánh bằng \texttt{timingSafeEqual}.
  \item \textbf{Phiên đăng nhập:} token \texttt{payload.chữ\_ký} ký HMAC-SHA256, hết hạn sau 12 giờ; middleware \texttt{requireStaff(...roles)} kiểm tra vai trò cho từng API.
  \item \textbf{SQL injection:} 100\% truy vấn dùng placeholder \texttt{?}; kiểm thử với chuỗi \verb|' OR 1=1 --| trả về danh sách rỗng.
  \item \textbf{XSS:} mọi dữ liệu người dùng (tên, tin nhắn chat) được escape trước khi chèn vào HTML.
  \item \textbf{Dữ liệu cá nhân:} tra cứu hồ sơ phải khớp cả mã và SĐT; sai một trong hai đều trả cùng thông báo 404 để không lộ thông tin.
\end{itemize}

\section{Kết quả giao diện}
\subsection{Cổng học viên}
\fig[0.92\linewidth]{01_trang_chu}{Trang chủ cổng học viên: quy định mở lớp và quy trình 4 bước}{fig:s1}
\fig[0.92\linewidth]{02_khoa_hoc}{Chức năng 1 – Quảng bá: danh mục khoá học, lọc theo chứng chỉ}{fig:s2}
\fig[0.92\linewidth]{03_chi_tiet_khoa_hoc}{Chức năng 1 – Chi tiết khoá học: lộ trình, chuẩn đầu ra và các lớp đang mở}{fig:s3}
\fig[0.92\linewidth]{04_tra_cuu_lop}{Chức năng 3 – Tra cứu lớp: lọc Offline + ca tối; thanh sĩ số có vạch ngưỡng mở lớp}{fig:s4}
\fig[0.92\linewidth]{05_dang_ky_validate}{Chức năng 2 – Form đăng ký với kiểm tra hợp lệ theo từng trường}{fig:s5}
\fig[0.92\linewidth]{18_tra_cuu_ho_so_thanh_toan}{Tra cứu hồ sơ: tiến trình 4 bước và nộp học phí}{fig:s18}
\fig[0.92\linewidth]{20_bai_kiem_tra}{Bài kiểm tra đầu vào: 20 câu, đồng hồ đếm ngược}{fig:s20}
\fig[0.92\linewidth]{21_tra_cuu_ho_so_da_thi}{Hồ sơ sau khi thi: kết quả CEFR đã chuyển về văn phòng}{fig:s21}
\fig[0.92\linewidth]{16_chat_hoc_vien_nhan_tin}{Chức năng 4 – Widget chat phía học viên nhận phản hồi realtime}{fig:s16}

\subsection{Dashboard quản trị}
\fig[0.92\linewidth]{08_dashboard}{Dashboard tổng quan: KPI, doanh thu – đăng ký 6 tháng, phễu tuyển sinh}{fig:s8}
\fig[0.92\linewidth]{09_dashboard_2}{Dashboard: sĩ số các lớp đang tuyển sinh so với ngưỡng, phân bố trình độ CEFR}{fig:s9}
\fig[0.92\linewidth]{11_xep_lop}{Xếp lớp: gợi ý lớp theo BR-06; lớp đủ 30 hiển thị nút “Lập tờ trình”}{fig:s11}
\fig[0.92\linewidth]{13_chi_tiet_lop}{Chi tiết lớp: lý do chưa đủ điều kiện mở, lập lịch theo thứ – ca, phân phòng}{fig:s13}
\fig[0.92\linewidth]{14_phe_duyet}{Phê duyệt sĩ số: hệ thống đề xuất “giữ ngoại lệ” (33 \texorpdfstring{$\leq$}{<=} 40 chỗ) và “chia đều 21 + 21” (42 học viên)}{fig:s14}
\fig[0.92\linewidth]{15_console_tu_van}{Chức năng 4 – Console tư vấn: hàng đợi phiên chat và hội thoại realtime}{fig:s15}
\fig[0.4\linewidth]{17_mobile_dark}{Giao diện trên điện thoại ở chế độ tối}{fig:s17}

\chapter{Kiểm thử}
\label{ch:test}

\section{Chiến lược kiểm thử}
Hệ thống được kiểm thử tự động ở ba mức, chạy bằng một lệnh \texttt{npm test}:
\begin{itemize}
  \item \textbf{Kiểm thử đơn vị} (\texttt{tests/rules.test.js}, 39 ca): kiểm tra các hàm thuần của bộ quy tắc nghiệp vụ, tập trung vào \emph{giá trị biên}: 9/10 học viên (online), 19/20 và 30/31 (offline), 39/40 (chia đều), các mốc quy đổi điểm CEFR.
  \item \textbf{Kiểm thử tích hợp API} (\texttt{tests/api.test.js}, 31 ca): khởi động máy chủ thật trên cổng ngẫu nhiên với \emph{CSDL tạm riêng} được nạp dữ liệu mẫu, gọi HTTP như trình duyệt, đi hết luồng đăng ký → thanh toán → thi → xếp lớp → tờ trình → phê duyệt → lập lịch → mở/huỷ lớp, đồng thời kiểm tra phân quyền.
  \item \textbf{Kiểm thử thời gian thực} (\texttt{tests/chat.test.js}, 8 ca): các client Socket.IO (một khách, hai tư vấn viên, một giáo vụ) kiểm tra hàng đợi, nhận phiên tranh chấp, nhắn tin hai chiều, khôi phục phiên, chặn socket lạ, đóng phiên.
\end{itemize}
Mỗi tệp kiểm thử dùng CSDL riêng trong thư mục tạm nên kết quả lặp lại được và không ảnh hưởng dữ liệu demo.

\section{Các ca kiểm thử tiêu biểu}
\begin{xltabular}{\linewidth}{|l|L{5.4cm}|Y|c|}
  \caption{Ca kiểm thử tiêu biểu và kết quả}\label{tab:tc}\\ \hline
  \headrow \hd{Mã} & \hd{Kịch bản} & \hd{Kết quả mong đợi} & \hd{KQ} \\ \hline
  \endfirsthead \hline
  \headrow \hd{Mã} & \hd{Kịch bản} & \hd{Kết quả mong đợi} & \hd{KQ} \\ \hline
  \endhead
  TC-01 & Lớp Online 9 / 10 học viên & BELOW\_MIN (thiếu 1) / ELIGIBLE & Đạt \\ \hline
  TC-02 & Lớp Offline 19, 20, 30, 31 học viên & BELOW\_MIN, ELIGIBLE, FULL, OVER\_CAPACITY & Đạt \\ \hline
  TC-03 & Xếp học viên thứ 30 / thứ 31 vào lớp Offline & ENROLLED / NEEDS\_APPROVAL & Đạt \\ \hline
  TC-04 & Xếp vào lớp Online đã đủ 40 & REJECT – ONLINE\_FULL, không lập tờ trình & Đạt \\ \hline
  TC-05 & Tờ trình 30 + 3, phòng 40 chỗ & Đề xuất KEEP; SPLIT dạng REMAINDER [30, 3] & Đạt \\ \hline
  TC-06 & Tờ trình 30 + 3, phòng 30 chỗ & KEEP không khả thi, đề xuất SPLIT & Đạt \\ \hline
  TC-07 & Tờ trình 30 + 10 & SPLIT dạng EVEN [20, 20] & Đạt \\ \hline
  TC-08 & Trùng phòng / trùng giảng viên cùng thứ – ca – thời gian & Phát hiện xung đột ROOM / TEACHER & Đạt \\ \hline
  TC-09 & Đăng ký hợp lệ & 201, mã \texttt{DK-YYYY-NNNN}, PENDING\_PAYMENT & Đạt \\ \hline
  TC-10 & Đăng ký trùng khoá; dữ liệu sai & 409; 400 kèm lỗi từng trường & Đạt \\ \hline
  TC-11 & Tra cứu hồ sơ đúng mã, sai SĐT & 404 & Đạt \\ \hline
  TC-12 & Làm bài khi chưa nộp học phí & 403 PAYMENT\_REQUIRED (BR-01) & Đạt \\ \hline
  TC-13 & Thanh toán 2 lần & Lần 2 trả 409 & Đạt \\ \hline
  TC-14 & Làm bài: trả lời đúng 15/20 & 20 câu không lộ đáp án; điểm 15, trình độ B2; hồ sơ TESTED & Đạt \\ \hline
  TC-15 & Gọi API quản trị không token; tư vấn viên xếp lớp & 401; 403 & Đạt \\ \hline
  TC-16 & Tham số tra cứu \verb|' OR 1=1 --| & 200, danh sách rỗng, không lỗi SQL & Đạt \\ \hline
  TC-17 & Giáo vụ ra quyết định tờ trình & 403 – chỉ Lãnh đạo & Đạt \\ \hline
  TC-18 & Lãnh đạo giữ ngoại lệ (hội trường 45 chỗ) & Lớp có 31 học viên, không còn hồ sơ chờ & Đạt \\ \hline
  TC-19 & Lãnh đạo tách lớp 30 + 12 = 42 & Hai lớp 21 + 21; lớp mới có \texttt{parent\_class\_id}, được gợi ý phòng trống & Đạt \\ \hline
  TC-20 & Từ chối tờ trình & Hồ sơ chờ trở về TESTED & Đạt \\ \hline
  TC-21 & Mở lớp Offline 21 HV chưa có phòng & 422 – “chưa được phân phòng học” & Đạt \\ \hline
  TC-22 & Gán phòng đang bận cùng ca & 422 SCHEDULE\_CONFLICT & Đạt \\ \hline
  TC-23 & Gán phòng trống rồi mở lớp & IN\_PROGRESS & Đạt \\ \hline
  TC-24 & Huỷ lớp & Toàn bộ học viên trở về danh sách chờ xếp lớp & Đạt \\ \hline
  TC-25 & Hai tư vấn viên cùng nhận một phiên & Người thứ hai bị từ chối & Đạt \\ \hline
  TC-26 & Tin nhắn hai chiều realtime & Bên kia nhận đúng nội dung, đúng người gửi & Đạt \\ \hline
  TC-27 & Socket lạ gửi tin vào phiên người khác & Bị từ chối & Đạt \\ \hline
  TC-28 & Đóng phiên rồi gửi tiếp & Bị từ chối; tin nhắn đã lưu CSDL & Đạt \\ \hline
  TC-29 & Tách lớp khi 30 + 35 hồ sơ chờ & Ba lớp 22 / 22 / 21 (không lớp nào vượt 30); hai lớp mới khác giảng viên, không trùng phòng & Đạt \\ \hline
  TC-30 & Sai mật khẩu 8 lần liên tiếp & Lần thứ 9 trả 429 kể cả khi nhập đúng & Đạt \\ \hline
  TC-31 & Giáo vụ nhận phiên chat & Bị từ chối (chỉ tư vấn viên/lãnh đạo) & Đạt \\ \hline
\end{xltabular}

\section{Kết quả chạy kiểm thử}
\begin{figure}[H]
\begin{PlainCode}
> node --test --test-reporter=spec "tests/*.test.js"
> BR-04 – phương án xử lý tờ trình
  v 30 + 3 = 33 < 40, phòng 40 chỗ -> đề xuất giữ ngoại lệ, tách phần dư
  v 30 + 3 nhưng phòng chỉ 30 chỗ -> không giữ được, đề xuất tách
  v 30 + 10 = 40 -> chia đều 20/20 (cả 2 lớp đủ điều kiện mở)
  ...
> Xếp lớp & tờ trình vượt sĩ số (BR-04)
  v Học viên thứ 30 vào lớp offline -> ENROLLED; thứ 31 -> chờ duyệt + tờ trình
  v Lãnh đạo "tách lớp" 30 + 12 = 42 -> chia đều 21 / 21, lớp mới có lịch + phòng
  ...
i tests 78
i suites 13
i pass 78
i fail 0
\end{PlainCode}
\captionof{lstlisting}{Kết quả \texttt{npm test} (rút gọn)}\label{lst:testout}
\end{figure}

\begin{table}[H]\centering
\caption{Tổng hợp kết quả kiểm thử tự động}\label{tab:testsum}
\begin{tabular}{|l|c|c|c|}
  \hline
  \headrow \hd{Tệp kiểm thử} & \hd{Số ca} & \hd{Đạt} & \hd{Thời gian} \\ \hline
  \texttt{rules.test.js} – quy tắc nghiệp vụ & 39 & 39 & $\approx$ 0,1 s \\ \hline
  \texttt{api.test.js} – tích hợp REST API & 31 & 31 & $\approx$ 1,5 s \\ \hline
  \texttt{chat.test.js} – Socket.IO & 8 & 8 & $\approx$ 1,1 s \\ \hline
  \textbf{Tổng} & \textbf{78} & \textbf{78 (100\%)} & \\ \hline
\end{tabular}
\end{table}

\section{Ma trận truy vết yêu cầu}
Bảng~\ref{tab:trace} nối từng yêu cầu của đề bài tới use case, mô-đun cài đặt và ca kiểm thử tương ứng, bảo đảm không yêu cầu nào bị bỏ sót và mọi quy tắc nghiệp vụ đều có kiểm thử.

\begin{xltabular}{\linewidth}{|L{2.7cm}|L{1.5cm}|Y|L{3.4cm}|}
  \caption{Truy vết yêu cầu – use case – mã nguồn – kiểm thử}\label{tab:trace}\\ \hline
  \headrow \hd{Yêu cầu} & \hd{Use case} & \hd{Mô-đun cài đặt} & \hd{Kiểm thử} \\ \hline
  \endfirsthead \hline
  \headrow \hd{Yêu cầu} & \hd{Use case} & \hd{Mô-đun cài đặt} & \hd{Kiểm thử} \\ \hline
  \endhead
  FR-01, FR-02 (CN1) & UC-01 & \texttt{services/catalog.js}, \texttt{public/js/portal.js} & API danh mục khoá học \\ \hline
  FR-03 (CN3) & UC-02 & \texttt{catalog.searchClasses} & API lọc lớp, TC-16 \\ \hline
  FR-04 (CN2) & UC-03 & \texttt{registration.create}, \texttt{rules.validateRegistration} & TC-09, TC-10 \\ \hline
  FR-05, FR-06 & UC-04 & \texttt{registration.pay}, \texttt{findByCodeAndPhone} & TC-11, TC-13 \\ \hline
  FR-07, BR-01 & UC-05 & \texttt{services/placement.js}, \texttt{rules.canTakePlacement} & TC-12, TC-14 \\ \hline
  FR-08 – FR-10 (CN4) & UC-06 & \texttt{sockets/chat.js}, \texttt{services/chat.js} & TC-25 – TC-28, TC-31 \\ \hline
  FR-11, BR-06 & UC-07 & \texttt{enrollment.enroll}, \texttt{rules.placementDecision} & TC-03, TC-04 \\ \hline
  FR-13, BR-04 & UC-08 & \texttt{enrollment.decide}, \texttt{rules.proposeResolution} & TC-05 – TC-07, TC-17 – TC-20, TC-29 \\ \hline
  FR-12, BR-05, BR-08 & UC-09 & \texttt{enrollment.updateArrangement}, \texttt{openClass}, \texttt{rules.findConflicts}, \texttt{canOpenClass} & TC-08, TC-21 – TC-24 \\ \hline
  BR-02, BR-03, BR-07 & – & \texttt{rules.evaluateClass}, view \texttt{v\_class\_enrollment} & TC-01, TC-02 \\ \hline
  FR-14 & UC-10 & \texttt{services/stats.js}, \texttt{public/js/admin.js} & API dashboard \\ \hline
  Bảo mật, phân quyền & – & \texttt{middleware/auth.js}, \texttt{middleware/throttle.js} & TC-15, TC-16, TC-30 \\ \hline
\end{xltabular}

\section{Kiểm thử giao diện}
Ngoài kiểm thử tự động, toàn bộ giao diện được chạy thử trên trình duyệt Chrome (độ phân giải 1440×900 và giả lập điện thoại 390×844) bằng công cụ \texttt{tools/screenshots.js}: công cụ tự đăng nhập từng vai trò, đi qua các màn hình, gõ tin nhắn chat ở cả hai phía và chụp ảnh – các ảnh trong Chương~\ref{ch:impl} được tạo ra từ kịch bản này. Quá trình chạy thử đã giúp phát hiện và sửa các lỗi: đường dẫn thư viện Chart.js bị chặn bởi trường \texttt{exports} của gói; xung đột tên lớp CSS giữa khối người dùng và bong bóng tin nhắn; lớp đang học bị gợi ý ưu tiên cho học viên mới (đã bổ sung trọng số giảm cho lớp đã khai giảng).

\chapter{Kết luận và hướng phát triển}
\label{ch:conclusion}

\section{Kết quả đạt được}
\begin{itemize}
  \item \textbf{Phân tích:} bóc tách đề bài thành 14 yêu cầu chức năng và 8 quy tắc nghiệp vụ có mã truy vết; làm rõ các điểm đề bài chưa quy định (sĩ số tối đa lớp online, ý nghĩa của “tách lớp”) bằng giả định tường minh.
  \item \textbf{Thiết kế:} 21 biểu đồ UML/ERD dạng mã nguồn (6 use case, 4 hoạt động, 4 tuần tự, 3 trạng thái, 2 lớp, 1 ERD, 1 kiến trúc); CSDL 14 bảng chuẩn 3NF với ràng buộc nghiệp vụ được cưỡng chế ngay trong CSDL.
  \item \textbf{Cài đặt:} ứng dụng web hoàn chỉnh gồm cổng học viên (4 chức năng bắt buộc + thanh toán + kiểm tra đầu vào) và dashboard văn phòng cho 3 vai trò; chat thời gian thực có hàng đợi, khôi phục phiên, chỉ báo đang soạn tin; dashboard thống kê trực quan.
  \item \textbf{Chất lượng:} 78 ca kiểm thử tự động đạt 100\%, bao phủ toàn bộ quy tắc sĩ số ở các giá trị biên và các luồng nghiệp vụ đầu – cuối.
  \item \textbf{Phát hiện trong phân tích:} chia đôi lớp chỉ hợp lệ khi tổng sĩ số $\geq 40$; hệ thống tự đề xuất phương án phù hợp (chia đều, tách phần dư, giữ ngoại lệ theo sức chứa phòng) giúp Lãnh đạo quyết định nhanh và đúng.
\end{itemize}

\section{Hạn chế}
\begin{itemize}
  \item Cổng thanh toán và việc gửi thông báo (email/SMS) mới ở mức mô phỏng.
  \item Chưa có chức năng điểm danh, nhập điểm từng buổi và đánh giá cuối khoá để “lưu trữ toàn bộ quá trình đào tạo” ở mức chi tiết.
  \item Ngân hàng câu hỏi kiểm tra đầu vào còn nhỏ (45 câu) và mới có dạng trắc nghiệm; chưa kiểm tra kỹ năng nghe, nói, viết.
  \item SQLite phù hợp quy mô một trung tâm; khi mở rộng nhiều cơ sở cần chuyển sang hệ quản trị CSDL máy chủ. Bộ đếm chống dò mật khẩu lưu trong bộ nhớ nên chỉ đúng với một tiến trình máy chủ.
  \item Token đăng nhập lưu ở \texttt{localStorage} của trình duyệt; bản triển khai thật nên dùng cookie \texttt{HttpOnly} và HTTPS.
\end{itemize}

\section{Hướng phát triển}
\begin{itemize}
  \item Tích hợp cổng thanh toán thật (VNPay/MoMo sandbox) với xác nhận qua IPN, gửi email/Zalo thông báo kết quả thi và lịch học.
  \item Bổ sung phân hệ đào tạo: điểm danh theo buổi, nhập điểm giữa kỳ – cuối kỳ, cảnh báo học viên vắng nhiều.
  \item Bài kiểm tra thích ứng (adaptive test) chọn câu theo năng lực, bổ sung phần nghe có audio.
  \item Hỗ trợ chatbot trả lời tự động các câu hỏi thường gặp ngoài giờ làm việc, chuyển sang tư vấn viên khi cần.
  \item Triển khai bằng Docker, chuyển CSDL sang PostgreSQL; bổ sung nhật ký kiểm toán (audit log) cho các quyết định của Lãnh đạo.
\end{itemize}


% ------------------------------------------------------------- TÀI LIỆU THAM KHẢO
\clearpage
\phantomsection
\addcontentsline{toc}{chapter}{Tài liệu tham khảo}
\markboth{Tài liệu tham khảo}{}
\renewcommand{\bibname}{Tài liệu tham khảo}
\begin{thebibliography}{99}\setlength{\itemsep}{2pt}
  \bibitem{fowler} M. Fowler, \textit{UML Distilled: A Brief Guide to the Standard Object Modeling Language}, 3rd ed. Addison-Wesley, 2003.
  \bibitem{omg} Object Management Group, \textit{Unified Modeling Language (UML) Specification}, Version 2.5.1, 2017.
  \bibitem{cockburn} A. Cockburn, \textit{Writing Effective Use Cases}. Addison-Wesley, 2000.
  \bibitem{sommerville} I. Sommerville, \textit{Software Engineering}, 10th ed. Pearson, 2015.
  \bibitem{tt01} Bộ Giáo dục và Đào tạo, \textit{Thông tư số 01/2014/TT-BGDĐT ban hành Khung năng lực ngoại ngữ 6 bậc dùng cho Việt Nam}, 24/01/2014.
  \bibitem{cefr} Council of Europe, \textit{Common European Framework of Reference for Languages: Learning, Teaching, Assessment – Companion Volume}. Strasbourg, 2020.
  \bibitem{express} OpenJS Foundation, \textit{Express 5.x – API Reference}. \url{https://expressjs.com}
  \bibitem{socketio} Socket.IO, \textit{Socket.IO v4 Documentation}. \url{https://socket.io/docs/v4/}
  \bibitem{sqlite} SQLite Consortium, \textit{SQLite Documentation – Partial Indexes, Views, Foreign Keys}. \url{https://www.sqlite.org/docs.html}
  \bibitem{bsqlite} J. Wise, \textit{better-sqlite3 – API documentation}. \url{https://github.com/WiseLibs/better-sqlite3}
  \bibitem{mermaid} Mermaid, \textit{Mermaid – Diagramming and charting tool}. \url{https://mermaid.js.org}
  \bibitem{chartjs} Chart.js, \textit{Chart.js 4 Documentation}. \url{https://www.chartjs.org/docs/}
  \bibitem{nodetest} Node.js, \textit{Test runner (node:test)}. \url{https://nodejs.org/api/test.html}
\end{thebibliography}

% ------------------------------------------------------------- PHỤ LỤC
\appendix
\chapter{Hướng dẫn cài đặt và chạy thử}
\label{app:setup}

\section{Yêu cầu môi trường}
Node.js $\geq$ 22 (đã thử trên Node.js 24.11), trình duyệt Chrome hoặc Edge. Không cần cài đặt máy chủ CSDL. Để biên dịch báo cáo cần MiKTeX hoặc TeX Live có XeLaTeX.

\section{Các lệnh}
\begin{figure}[H]
\begin{PlainCode}
cd ProjectII
npm install            # cài thư viện
npm run seed           # (tuỳ chọn) tạo lại CSDL mẫu database.sqlite
npm start              # http://localhost:3000  (cổng học viên)
                       # http://localhost:3000/admin  (văn phòng)
npm test               # chạy 72 ca kiểm thử tự động
npm run diagrams       # render lại 21 biểu đồ Mermaid -> docs/diagrams/png
npm run screenshots    # chụp lại ảnh giao diện (khi máy chủ đang chạy)
cd report && latexmk -xelatex main.tex   # biên dịch báo cáo
\end{PlainCode}
\captionof{lstlisting}{Cài đặt và chạy}\label{lst:run}
\end{figure}
Trong VS Code (tiện ích LaTeX Workshop), chọn recipe \textit{latexmk (XeLaTeX)} đã được cấu hình sẵn trong \texttt{.vscode/settings.json}.

\section{Tài khoản và dữ liệu demo}
\begin{table}[H]\centering
\caption{Tài khoản nhân viên (mật khẩu \texttt{123456})}\label{tab:accounts}
\begin{tabular}{|l|l|l|}
  \hline
  \headrow \hd{Tên đăng nhập} & \hd{Họ tên} & \hd{Vai trò} \\ \hline
  \texttt{tuvan1}, \texttt{tuvan2} & Phạm Thu Trang, Lê Minh Đức & Tư vấn viên \\ \hline
  \texttt{giaovu} & Trần Thị Mai & Giáo vụ \\ \hline
  \texttt{lanhdao} & Nguyễn Văn Hùng & Lãnh đạo trung tâm \\ \hline
\end{tabular}
\end{table}
Dữ liệu mẫu gồm 8 khoá học, 20 lớp, 8 giảng viên, 7 phòng học, 410 hồ sơ, 45 câu hỏi kiểm tra và 4 phiên chat, được thiết kế để phủ mọi trạng thái nghiệp vụ. Ba học viên demo có SĐT cố định để thử tra cứu hồ sơ: \texttt{0912345678} (chờ nộp phí), \texttt{0987654321} (đã nộp – chờ thi), \texttt{0905123456} (đã thi – chờ xếp lớp); mã hồ sơ tương ứng được in ra khi chạy \texttt{npm run seed}.

\section{Kịch bản chạy thử đầu – cuối}
\begin{enumerate}
  \item \textbf{Cổng học viên:} xem khoá học (CN1) → lọc lớp Offline, ca tối (CN3) → bấm “Đăng ký lớp này” và gửi form (CN2) → “Nộp học phí ngay” → làm bài kiểm tra → xem kết quả CEFR.
  \item \textbf{Chat (CN4):} mở widget chat ở cổng học viên; ở cửa sổ khác đăng nhập \texttt{tuvan1}, vào “Tư vấn trực tuyến”, nhận phiên và trả lời – tin nhắn xuất hiện tức thì ở hai phía.
  \item \textbf{Xếp lớp:} đăng nhập \texttt{giaovu} → “Hồ sơ \& xếp lớp” → chọn hồ sơ vừa thi → xếp vào lớp \texttt{VSTEP-B1-B1-OFF-01} (đang 29/30): học viên thứ 30 được xếp ngay; xếp thêm một học viên nữa → hệ thống tự lập tờ trình.
  \item \textbf{Phê duyệt:} đăng nhập \texttt{lanhdao} → “Phê duyệt sĩ số”: tờ trình lớp \texttt{VSTEP-B1-A2-OFF-02} (30 + 12) được đề xuất chia đều 21 + 21; tờ trình lớp \texttt{IELTS-I-B1-OFF-02} (30 + 3, phòng 40 chỗ) được đề xuất giữ ngoại lệ.
  \item \textbf{Mở lớp:} lớp \texttt{VSTEP-B2-B2-OFF-01} có 21 học viên nhưng chưa có phòng → không mở được; phân phòng P.302 ca chiều T3–T5 → mở lớp thành công.
\end{enumerate}


\end{document}
